\documentclass[11pt]{article}
\usepackage[margin=1in]{geometry}
\usepackage{amsmath,amssymb}
\usepackage{booktabs}
\usepackage{graphicx}
\usepackage{hyperref}

\title{Off-Policy KTM Pipeline: End-to-End Methodology}
\author{Geometric Flow Repository}
\date{\today}

\begin{document}
\maketitle

\section{Goal}
This document describes the full experimental pipeline implemented in:
\begin{itemize}
  \item \texttt{offpolicy\_ktm\_pipeline/run\_full\_pipeline.py}
  \item \texttt{offpolicy\_ktm\_pipeline/src/*.py}
\end{itemize}

The pipeline has two main parts:
\begin{enumerate}
  \item Build a contextual logging dataset $(\theta,\delta,r,o)$ from student-item logs.
  \item Learn and evaluate a new policy for item difficulty selection with off-policy estimators.
\end{enumerate}

\section{Data Processing and KTM Latents}
\subsection{PIX preprocessing}
For PIX-like data, \texttt{data/pix\_processing.py} applies:
\begin{itemize}
  \item ID remapping to contiguous integers for user/item/skill.
  \item \texttt{answer\_result} mapping:
  \[
  \texttt{ok}\mapsto 1,\quad \texttt{ko}\mapsto 0,\quad \texttt{aband}\mapsto 0.5
  \]
  \item Binary outcome column:
  \[
  \texttt{outcome} =
  \begin{cases}
  0 & \text{if } \texttt{answer\_result}=0.5\\
  \texttt{answer\_result} & \text{otherwise}
  \end{cases}
  \]
\end{itemize}

\subsection{KTM-style latent construction}
\texttt{models/ktm.py::processing\_data} fits a one-hot user-item logistic model:
\[
\Pr(y=1\mid u,i)=\sigma\!\left(\beta_u+\gamma_i+b\right)
\]
Then defines:
\[
\theta_u=\beta_u,\qquad \delta_i=-\gamma_i
\]
where $\theta_u$ is proficiency and $\delta_i$ is item difficulty.

It then:
\begin{itemize}
  \item ranks users and items by latent value (for stable integer IDs),
  \item attaches \texttt{proficiency} and \texttt{difficulties},
  \item builds sequence fields:
  \[
  \texttt{sequence\_position},\ \texttt{sequence\_length},\ \texttt{order\_sequence}.
  \]
\end{itemize}

\section{Round Construction}
The pipeline trains/evaluates by rounds indexed by \texttt{order\_sequence}. If needed, long tails are compressed to a fixed number of rounds (\texttt{--max-rounds}) with:
\begin{itemize}
  \item \texttt{tail}: keep early rounds, merge tail into final bucket.
  \item \texttt{quantile}: weighted quantile bucketing by observation count.
  \item \texttt{tail\_count\_at\_n}: after round $N$, aggregate tail windows to roughly match count at round $N$.
\end{itemize}

\section{Behavior Policy via Gaussian Regression}
For each round $o$, \texttt{fit\_gaussian\_by\_order} fits:
\[
\delta \mid \theta,o \sim \mathcal{N}(\mu_o(\theta),\sigma_o(\theta)^2)
\]
with polynomial parameterization:
\[
\mu_o(\theta)=\sum_{d=0}^{D_\mu}\beta^{(\mu)}_{o,d}\theta^d,\qquad
\log \sigma_o(\theta)=\sum_{d=0}^{D_\sigma}\beta^{(\sigma)}_{o,d}\theta^d
\]
and hard floor:
\[
\sigma_o(\theta)=\max\!\left(\exp(\log \sigma_o(\theta)),\sigma_{\min}\right),
\]
where \texttt{sigma\_floor} is typically $0.2$.

The per-round fit minimizes average Gaussian negative log-likelihood:
\[
\mathcal{L}_{\text{NLL}}=
\frac{1}{n}\sum_{j=1}^n
\left[
\frac{1}{2}\log(2\pi)+\log \sigma_o(\theta_j)
+\frac{1}{2}\left(\frac{\delta_j-\mu_o(\theta_j)}{\sigma_o(\theta_j)}\right)^2
\right].
\]

\section{Propensity and Reward}
After attaching round-specific behavior parameters, each tuple gets:
\[
\log \pi_b(\delta\mid\theta,o)=
-\frac{1}{2}\log(2\pi)-\log \sigma_o(\theta)
-\frac{1}{2}\left(\frac{\delta-\mu_o(\theta)}{\sigma_o(\theta)}\right)^2.
\]

Reward is defined as:
\[
r = \texttt{correct}\cdot(\delta-\delta_{\min}),
\]
so successful responses on harder items get larger reward.

\section{Learned Policy and Training Objective}
\subsection{Policy class}
\texttt{models/offpolicy\_gaussian\_policy.py::GlobalGaussianPolicy} learns one global policy:
\[
\pi_\phi(\delta\mid\theta)=\mathcal{N}(\mu_\phi(\theta),\sigma_\phi(\theta)^2),
\]
\[
\mu_\phi(\theta)=\sum_{d=0}^{D_\mu}\alpha_d\theta^d,\qquad
\log \sigma_\phi(\theta)=\sum_{d=0}^{D_\sigma}\eta_d\theta^d,\qquad
\sigma_\phi(\theta)\ge \sigma_{\min}.
\]

\subsection{Importance weights}
For each sample:
\[
w=\exp\left(\log \pi_\phi(\delta\mid\theta)-\log \pi_b(\delta\mid\theta,o)\right),
\qquad
\tilde{w}=\min(w,c_{\text{clip}}).
\]
The code also clips log-ratio and max weight for stability.

\subsection{Optimized objectives}
Supported objectives (\texttt{--objective}):
\[
\widehat{V}_{\text{IPS}}=\frac{1}{n}\sum_i w_i r_i,\quad
\widehat{V}_{\text{SNIPS}}=\frac{\sum_i w_i r_i}{\sum_i w_i},
\]
\[
\widehat{V}_{\text{CIPS}}=\frac{1}{n}\sum_i \tilde{w}_i r_i,\quad
\widehat{V}_{\text{CSNIPS}}=\frac{\sum_i \tilde{w}_i r_i}{\sum_i \tilde{w}_i}.
\]

Training minimizes:
\[
\mathcal{J}(\phi)=
-\widehat{V}_{\text{obj}}
+\lambda_2\left(\|\alpha\|_2^2+\|\eta\|_2^2\right)
+\lambda_{\text{var}}\sqrt{\frac{1}{t}}\sqrt{\operatorname{Var}(\text{contrib})},
\]
where $t$ is minibatch size and \texttt{contrib} is the per-sample term used by the selected estimator.

\section{Context Reweighting}
To handle context shift across rounds, the pipeline can use a target context distribution $q(\theta)$ (for example \texttt{best\_round\_in\_k}) and per-sample ratio:
\[
\omega(\theta)=\frac{q(\theta)}{p_t(\theta)}
\]
estimated from histograms and clipped. These context weights are used:
\begin{itemize}
  \item in training (\texttt{--context-match-sampling}),
  \item in evaluation (reported \texttt{*\_q\_*} metrics).
\end{itemize}

\section{Model-Based IRT Value Metrics}
The pipeline reports model-based values using:
\[
\Pr(\text{correct}=1\mid\theta,\delta)=\sigma(\theta-\delta),
\quad
r(\theta,\delta)=\sigma(\theta-\delta)(\delta-\delta_{\min}).
\]

For policy $\pi$:
\[
V_{\text{IRT}}(\pi)=\mathbb{E}_{\theta\sim p}\mathbb{E}_{\delta\sim\pi(\cdot\mid\theta)}[r(\theta,\delta)],
\]
estimated by Monte Carlo (\texttt{mc\_policy\_value}).

The oracle curve is:
\[
\delta^\star(\theta)\in\arg\max_{\delta\in[\delta_{\min},\delta_{\max}]}
\sigma(\theta-\delta)(\delta-\delta_{\min}),
\]
and the corresponding value is reported as \texttt{optimal\_irt\_value}.

\section{Experiment}
\subsection{Protocol}
We ran the simplified cumulative-last-round pipeline in
\texttt{offpolicy\_ktm\_pipeline/run\_last\_round\_simple.py} on three datasets:
\begin{itemize}
  \item PIX (\texttt{pix\_processed.csv}),
  \item SkillBuilder (\texttt{skill\_builder\_2015\_mintrace10.csv}),
  \item Assistments-like (\texttt{data\_with\_answer\_number.csv}).
\end{itemize}

Common settings:
\begin{itemize}
  \item \texttt{ktm\_c=0.1}, strict 1PL KTM with no intercept;
  \item behavior model: truncated Gaussian with $\mu(\theta)$ degree 1 and $\log\sigma(\theta)$ degree 2;
  \item context target: ``best round in first $N/2$'' (closest to last-round context);
  \item context reweighting with histogram ratio, \texttt{context\_bins=60}, \texttt{context\_ratio\_clip=10};
  \item IS weight cap \texttt{max\_weight=20}, policy sigma floor \texttt{0.2};
  \item rounds capped to 20 (\texttt{tail} strategy), training scope cumulative at final round;
  \item objectives trained separately: IPS, SNIPS, DR, MIS, DM.
\end{itemize}

Batch size followed a rule-of-thumb targeting roughly $50$--$100$ updates per epoch:
\texttt{4096} (Assistments), \texttt{8192} (SkillBuilder), \texttt{32768} (PIX).

\subsection{Final-Round Learned Policy Values}
Table~\ref{tab:final_learned} reports the final-round values of the \emph{learned policy}
under each estimator. Two robust trends are visible:
\begin{itemize}
  \item For PIX, DR/SNIPS/MIS policies are close in value, with large ESS.
  \item For Assistments and SkillBuilder, IPS/MIS-optimized policies tend to give
  higher IPS/MIS metrics, while DR/SNIPS-optimized policies are more conservative.
\end{itemize}
This confirms that objective choice primarily changes the policy shape in regions with
strong weighting effects, while remaining in a similar performance band for stable datasets.

\begin{table}[htbp]
\centering
\small
\caption{Final-round learned policy values across datasets/objectives (ESS in thousands).}
\label{tab:final_learned}
\begin{tabular}{llrrrrrrr}
\toprule
Dataset & Obj & IPS & SNIPS & DR & MIS & DM & ESS$_{log}$ (k) & ESS$_{mix}$ (k) \\
\midrule
assistments & ips & 0.703 & 0.677 & 0.675 & 0.708 & 0.616 & 157.677 & 155.989 \\
assistments & snips & 0.698 & 0.676 & 0.676 & 0.704 & 0.621 & 157.512 & 155.994 \\
assistments & dr & 0.694 & 0.676 & 0.676 & 0.701 & 0.626 & 156.845 & 155.567 \\
assistments & mis & 0.706 & 0.679 & 0.678 & 0.712 & 0.616 & 155.006 & 153.381 \\
assistments & dm & 0.515 & 0.547 & 0.566 & 0.516 & 0.724 & 86.121 & 100.156 \\
skillbuilder & ips & 1.128 & 1.020 & 1.005 & 1.127 & 0.993 & 419.337 & 422.537 \\
skillbuilder & snips & 1.080 & 0.997 & 0.978 & 1.080 & 0.967 & 441.189 & 443.839 \\
skillbuilder & dr & 1.048 & 0.991 & 0.972 & 1.049 & 0.964 & 511.248 & 516.276 \\
skillbuilder & mis & 1.125 & 1.020 & 1.003 & 1.125 & 0.992 & 422.379 & 425.624 \\
skillbuilder & dm & 1.021 & 0.980 & 0.961 & 1.023 & 0.954 & 540.528 & 546.522 \\
pix & ips & 2.638 & 2.864 & 2.863 & 2.933 & 2.582 & 497.940 & 908.155 \\
pix & snips & 2.561 & 2.853 & 2.851 & 2.902 & 2.572 & 424.182 & 875.189 \\
pix & dr & 2.590 & 2.865 & 2.867 & 2.913 & 2.586 & 423.558 & 844.858 \\
pix & mis & 2.626 & 2.855 & 2.859 & 2.935 & 2.575 & 465.723 & 903.800 \\
pix & dm & 2.576 & 2.855 & 2.856 & 2.889 & 2.584 & 429.977 & 873.509 \\
\bottomrule
\end{tabular}

\end{table}

\subsection{Adaptive Behavior Policy by Round}
The adaptive behavior policy value was estimated at \emph{every} round using
context-reweighted estimators on cumulative seen data, and we also report
context-reweighted batch reward means. Full per-round values are included in
the generated tables below.

\begin{table}[p]
\centering
\scriptsize
\caption{Assistments: behavior policy values by round (context-reweighted).}
\label{tab:beh_round_assistments}
\begin{tabular}{rrrrrrr}
\toprule
Round & BatchReward$_q$ & IPS$_q$ & SNIPS$_q$ & DR$_q$ & MIS$_q$ & DM$_q$ \\
\midrule
1 & 0.563 & 0.554 & 0.552 & 0.552 & 0.554 & 0.634 \\
2 & 0.579 & 0.562 & 0.550 & 0.544 & 0.556 & 0.631 \\
3 & 0.600 & 0.570 & 0.567 & 0.567 & 0.567 & 0.637 \\
4 & 0.574 & 0.565 & 0.550 & 0.543 & 0.559 & 0.641 \\
5 & 0.619 & 0.575 & 0.567 & 0.565 & 0.571 & 0.638 \\
6 & 0.599 & 0.576 & 0.569 & 0.569 & 0.573 & 0.641 \\
7 & 0.636 & 0.580 & 0.568 & 0.565 & 0.578 & 0.641 \\
8 & 0.613 & 0.584 & 0.579 & 0.578 & 0.583 & 0.643 \\
9 & 0.635 & 0.593 & 0.591 & 0.592 & 0.592 & 0.639 \\
10 & 0.631 & 0.592 & 0.589 & 0.588 & 0.591 & 0.644 \\
11 & 0.622 & 0.591 & 0.581 & 0.575 & 0.589 & 0.643 \\
12 & 0.631 & 0.591 & 0.581 & 0.579 & 0.590 & 0.646 \\
13 & 0.612 & 0.593 & 0.584 & 0.579 & 0.592 & 0.649 \\
14 & 0.610 & 0.591 & 0.584 & 0.582 & 0.591 & 0.651 \\
15 & 0.632 & 0.593 & 0.585 & 0.584 & 0.593 & 0.646 \\
16 & 0.612 & 0.592 & 0.583 & 0.582 & 0.592 & 0.649 \\
17 & 0.615 & 0.593 & 0.583 & 0.578 & 0.592 & 0.650 \\
18 & 0.632 & 0.591 & 0.582 & 0.579 & 0.590 & 0.650 \\
19 & 0.625 & 0.594 & 0.579 & 0.573 & 0.592 & 0.648 \\
20 & 0.625 & 0.608 & 0.603 & 0.604 & 0.608 & 0.652 \\
\bottomrule
\end{tabular}

\end{table}

\begin{table}[p]
\centering
\scriptsize
\caption{SkillBuilder: behavior policy values by round (context-reweighted).}
\label{tab:beh_round_skillbuilder}
\begin{tabular}{rrrrrrr}
\toprule
Round & BatchReward$_q$ & IPS$_q$ & SNIPS$_q$ & DR$_q$ & MIS$_q$ & DM$_q$ \\
\midrule
1 & 0.656 & 0.570 & 0.568 & 0.583 & 0.570 & 0.704 \\
2 & 0.763 & 0.609 & 0.606 & 0.621 & 0.608 & 0.699 \\
3 & 0.806 & 0.633 & 0.630 & 0.645 & 0.633 & 0.697 \\
4 & 0.803 & 0.652 & 0.653 & 0.664 & 0.651 & 0.708 \\
5 & 0.830 & 0.659 & 0.659 & 0.671 & 0.659 & 0.709 \\
6 & 0.842 & 0.665 & 0.664 & 0.677 & 0.665 & 0.708 \\
7 & 0.809 & 0.670 & 0.665 & 0.681 & 0.670 & 0.711 \\
8 & 0.819 & 0.672 & 0.666 & 0.683 & 0.672 & 0.711 \\
9 & 0.825 & 0.673 & 0.665 & 0.685 & 0.673 & 0.709 \\
10 & 0.822 & 0.680 & 0.673 & 0.693 & 0.680 & 0.715 \\
11 & 0.817 & 0.682 & 0.674 & 0.694 & 0.681 & 0.714 \\
12 & 0.825 & 0.686 & 0.677 & 0.698 & 0.685 & 0.716 \\
13 & 0.828 & 0.689 & 0.682 & 0.701 & 0.689 & 0.718 \\
14 & 0.824 & 0.693 & 0.685 & 0.704 & 0.692 & 0.721 \\
15 & 0.827 & 0.693 & 0.686 & 0.704 & 0.692 & 0.720 \\
16 & 0.835 & 0.695 & 0.689 & 0.707 & 0.695 & 0.722 \\
17 & 0.831 & 0.696 & 0.690 & 0.708 & 0.696 & 0.723 \\
18 & 0.846 & 0.697 & 0.691 & 0.709 & 0.697 & 0.724 \\
19 & 0.846 & 0.699 & 0.694 & 0.711 & 0.699 & 0.724 \\
20 & 0.902 & 0.741 & 0.737 & 0.756 & 0.744 & 0.753 \\
\bottomrule
\end{tabular}

\end{table}

\begin{table}[p]
\centering
\scriptsize
\caption{PIX: behavior policy values by round (context-reweighted).}
\label{tab:beh_round_pix}
\begin{tabular}{rrrrrrr}
\toprule
Round & BatchReward$_q$ & IPS$_q$ & SNIPS$_q$ & DR$_q$ & MIS$_q$ & DM$_q$ \\
\midrule
1 & 0.325 & 3.347 & 1.726 & 1.555 & 3.347 & 1.357 \\
2 & 1.074 & 1.855 & 1.309 & 0.902 & 0.777 & 0.806 \\
3 & 1.726 & 0.655 & 0.581 & 0.602 & 0.413 & 0.595 \\
4 & 2.009 & 0.763 & 0.746 & 0.719 & 0.492 & 0.709 \\
5 & 2.526 & 1.527 & 1.346 & 1.206 & 1.019 & 1.164 \\
6 & 2.683 & 1.873 & 1.717 & 1.666 & 1.667 & 1.602 \\
7 & 2.666 & 1.336 & 1.514 & 1.366 & 1.234 & 1.343 \\
8 & 2.757 & 1.833 & 1.896 & 1.853 & 1.909 & 1.787 \\
9 & 2.638 & 2.285 & 2.401 & 2.505 & 2.526 & 2.364 \\
10 & 2.608 & 2.117 & 2.719 & 2.764 & 2.792 & 2.590 \\
11 & 2.358 & 1.646 & 2.646 & 2.567 & 2.706 & 2.444 \\
12 & 2.349 & 1.753 & 2.668 & 2.606 & 2.709 & 2.473 \\
13 & 2.182 & 1.663 & 2.584 & 2.500 & 2.558 & 2.390 \\
14 & 1.789 & 1.066 & 2.025 & 1.686 & 1.665 & 1.727 \\
15 & 1.808 & 1.101 & 1.991 & 1.725 & 1.701 & 1.773 \\
16 & 2.353 & 2.368 & 2.703 & 2.691 & 2.702 & 2.483 \\
17 & 1.947 & 1.650 & 2.234 & 2.091 & 2.095 & 2.046 \\
18 & 2.038 & 1.696 & 2.251 & 2.125 & 2.131 & 2.070 \\
19 & 2.081 & 1.768 & 2.282 & 2.169 & 2.177 & 2.102 \\
20 & 2.074 & 1.985 & 2.281 & 2.248 & 2.249 & 2.205 \\
\bottomrule
\end{tabular}

\end{table}

\subsection{Policy Curves}
For each dataset we plot:
\begin{enumerate}
  \item Learned policies from all estimators as mean difficulty
  $\mathbb{E}[\delta\mid\theta]$ with the optimal IRT policy shown as dotted curve.
  \item The same learned curves plus behavior policy snapshots (round 1, mid, last).
\end{enumerate}

\begin{figure}[p]
\centering
\includegraphics[width=0.48\textwidth]{../../pix_mapping/multi_dataset_last_round_report/assistments_policy_curves_all_estimators.png}
\includegraphics[width=0.48\textwidth]{../../pix_mapping/multi_dataset_last_round_report/assistments_policy_curves_with_behavior_rounds.png}
\caption{Assistments policy curves: learned estimators vs optimal IRT, and behavior evolution.}
\label{fig:assistments_curves}
\end{figure}

\begin{figure}[p]
\centering
\includegraphics[width=0.48\textwidth]{../../pix_mapping/multi_dataset_last_round_report/skillbuilder_policy_curves_all_estimators.png}
\includegraphics[width=0.48\textwidth]{../../pix_mapping/multi_dataset_last_round_report/skillbuilder_policy_curves_with_behavior_rounds.png}
\caption{SkillBuilder policy curves: learned estimators vs optimal IRT, and behavior evolution.}
\label{fig:skillbuilder_curves}
\end{figure}

\begin{figure}[p]
\centering
\includegraphics[width=0.48\textwidth]{../../pix_mapping/multi_dataset_last_round_report/pix_policy_curves_all_estimators.png}
\includegraphics[width=0.48\textwidth]{../../pix_mapping/multi_dataset_last_round_report/pix_policy_curves_with_behavior_rounds.png}
\caption{PIX policy curves: learned estimators vs optimal IRT, and behavior evolution.}
\label{fig:pix_curves}
\end{figure}

\subsection{Discussion}
Across datasets, the objective family mostly changes policy aggressiveness rather than
completely reordering performance. In PIX, all non-DM objectives remain close with high ESS,
suggesting low variance pressure after clipping. In Assistments/SkillBuilder, differences are
larger, and DR/SNIPS provide more stable trade-offs against ESS than raw IPS-like objectives.
Behavior-evolution plots show that learned policies generally move away from early-round behavior
toward higher expected difficulty in high-$\theta$ regions, while the dotted IRT-optimal policy
remains an upper reference envelope.

\section{Main Outputs}
Each run writes:
\begin{itemize}
  \item \texttt{ktm\_dataframe.csv}
  \item \texttt{round\_bucket\_mapping.csv}
  \item \texttt{gaussian\_fit\_by\_order.csv}
  \item \texttt{offpolicy\_tuples.csv}
  \item \texttt{policy\_round\_results.csv}
  \item \texttt{policy\_value\_table.csv}
  \item \texttt{policy\_coefficients\_by\_round.csv}
  \item optional plots/animation.
\end{itemize}

\section{Reference Command}
\begin{verbatim}
python offpolicy_ktm_pipeline/run_full_pipeline.py \
  --input-csv pix_mapping/pix_processed.csv \
  --objective ips \
  --denominator mixture \
  --train-scope cumulative \
  --context-target best_round_in_k \
  --context-k-rounds 5 \
  --balanced-minibatch \
  --context-match-sampling \
  --sigma-floor 0.2
\end{verbatim}

\end{document}
